% NOETHER PAPER 4 — KOREAN TRANSLATION-PRODUCER DRAFT — UNCHECKED
% Unit: T05-U22; current-authority whole lines 4003--4016
% Source slice: 956 LF UTF-8 bytes; SHA-256 A9DCEA57176D7C441AB85678D6ED340BA2C65C7C9596C4B2B650C92C780897D5
% Pointer: NOETH-DE-AUTH-v004-20260804; SHA-256 A1C62FDACAA34DFC1B806DC18258F2E732539F7AE9D85AA4BD9E1067B8749D9F
% This is producer translation only: no source/Korean/formula review, compilation, or rendering.

\section*{\S\ 5. 명시적 형태의 분해 항등식.}
경우 4)의 (25.)에서 $\tau$를 $\sigma+\rho$로 특수화하면 다음이 나온다.
\begin{equation}
 \pair{q_\rho A^\sigma}{y^{(1)}\cdots y^{(\rho+\sigma)}}
 =\sum_{i_\rho=i_{\rho-1}+1}^{\tau}
 \eps_{i_\rho}\pair{q_\rho}{y^{(i_1)}\cdots y^{(i_\rho)}}
 \pair{A^\sigma}{y^{(i_{\rho+1})}\cdots y^{(i_{\rho+\sigma})}},
\tag{32}
\end{equation}
여기서 $\eps_{i_\rho}=(-1)^{\delta_{i_\rho}}$, $\delta_{i_\rho}=\rho(\rho+1)/2+i_1+\cdots+i_\rho$가 된다. 이것은 (16.)과 (17.)보다 더 일반적인 형태의 행렬 곱정리이며, 곱행렬의 행렬식
\[
 \sum\pm(v^{(1)}y^{(1)})\cdots(v^{(\rho)}y^{(\rho)})(a^{(1)}y^{(\rho+1)})\cdots(a^{(\sigma)}y^{(\rho+\sigma)})
\]
을 라플라스 전개하여 얻는다. 따라서 더 일반적인 곱정리는 항등식 (20.), 또는 각각 (21.)로부터 합성되며, 이로써 특수한 형태 (16.), (17.)을 따로 골라낸 이유가 설명된다.
